\chapter{Users and Output Guide}
\label{chap:users_outputs}

\section{Who Uses AnsibleGuard?}
AnsibleGuard is not designed for one generic ``user.'' Different actors need
different evidence, latency, and detail. The same canonical model supports all
of them without forcing each actor to interpret raw provider output.

\begin{longtable}{p{3.1cm}p{3.5cm}p{3.2cm}p{4.5cm}}
\toprule User type & Primary question & Recommended mode/profile & Principal outputs \\
\midrule\endhead
Ansible author & What should I correct before committing? & Local path or
changed-file feedback; \code{developer-fast} or \code{quality} & Text/Markdown
for explanation; per-file findings and exact source lines. \\
Role or collection maintainer & Is the complete role structurally and
semantically sound? & Full repository; \code{technical-full}, \code{quality},
or \code{security} & Canonical occurrences, category/rule summary, role-layout,
Molecule, documentation and dependency findings. \\
Pull-request reviewer & What did this PR introduce or fix? & Base/head PR mode;
\code{pr} plus \code{review-range} & Introduced, fixed, unchanged and
moved/uncertain sets,
review metadata, Markdown summary and SARIF annotations. \\
CI/CD engineer & Should this pipeline pass, warn, or fail? & Frozen commit with
a project profile/policy & SARIF, JSON contract, analysis health, severity and
diagnostic fields; stable exit behavior. \\
Security engineer & Are credentials, unsafe permissions or insecure transport
present? & \code{security} with GitLeaks and selected ansible-lint rules &
Redacted secret evidence, provider fingerprints, security concepts and
EXPOSED/uncertainty context. \\
Quality lead / architect & Which mechanisms dominate maintenance work? & Full
or \code{quality} profile across repositories & Category/concept distributions,
density denominators, structural and cross-file evidence. \\
Researcher & Can the study be reproduced and compared without population
drift? & Manifest-required snapshot profiles & Frozen SHA, manifest hash,
source/analysis population IDs, raw and canonical JSON/CSV, diagnostics and
schema copy. \\
Tool developer & Why did a rule or provider behave this way? & Focused tests,
single repository, raw + canonical outputs & Raw provider records, parsed-file
cache, metrics, diagnostics, native IDs, fingerprints and canonical evidence. \\
Supervisor / auditor & Which claims are supported by which artifacts? & Frozen
study aggregation and report & All-repositories CSV, methods, charts,
provider-health evidence, version comparison and stated validity limits. \\
\bottomrule
\end{longtable}

\section{User-to-Decision Flow}
\begin{figure}[H]
\centering
\resizebox{\textwidth}{!}{%
\begin{tikzpicture}[
 u/.style={draw=ansblue,rounded corners,fill=anslightblue,text width=27mm,
 minimum height=10mm,align=center,font=\small\bfseries},
 o/.style={draw=ansgreen,rounded corners,fill=anslightgreen,text width=30mm,
 minimum height=10mm,align=center,font=\small},
 a/.style={-{Latex[length=2mm]},thick},node distance=8mm and 12mm]
\node[u] (author) {Author / maintainer};
\node[u,below=of author] (reviewer) {Reviewer / CI};
\node[u,below=of reviewer] (researcher) {Researcher / auditor};
\node[o,right=of author] (fast) {Text, Markdown, per-file evidence};
\node[o,right=of reviewer] (diff) {SARIF and introduced/fixed comparison};
\node[o,right=of researcher] (data) {Validated JSON/CSV, raw evidence, metadata};
\node[o,right=20mm of diff] (decision) {Fix, approve, gate, investigate, or reproduce};
\draw[a] (author)--(fast); \draw[a] (reviewer)--(diff);
\draw[a] (researcher)--(data); \draw[a] (fast)--(decision);
\draw[a] (diff)--(decision); \draw[a] (data)--(decision);
\end{tikzpicture}}
\caption{Different users consume different projections of one evidence model.}
\end{figure}

\section{Actual Output Tree}
The following tree reflects the v0.0.6 technical-full snapshots. Compatibility
files remain available at the root while the structured directories separate
raw, canonical, metadata, diagnostics and schema responsibilities.

\begin{lstlisting}[basicstyle=\ttfamily\footnotesize]
<output>/
|-- contract.json
|-- diagnostics.json
|-- file_selection.json
|-- findings.json              # normalized compatibility projection
|-- findings.raw.json          # raw compatibility projection
|-- findings.csv
|-- findings.sarif
|-- findings.txt
|-- occurrences.json
|-- gitleaks-suppressed.json
|-- raw/
|   |-- findings.json
|   |-- findings.csv
|   `-- findings.txt
|-- canonical/
|   |-- occurrences.json
|   |-- occurrences.csv
|   |-- occurrences.txt
|   `-- findings.sarif
|-- metadata/
|   |-- run.json
|   |-- file-selection.json
|   `-- review-metrics.json
|-- files/<repository-relative-path>/
|   |-- metrics.json
|   |-- findings.json
|   |-- occurrences.json
|   |-- diagnostics.json
|   `-- report.json
|-- schemas/output-v3.json
|-- legacy/layout.json
`-- .cache/parsed/<content-hash>.json
\end{lstlisting}

The cache is an execution optimisation, not a study result. The legacy layout
exists for compatibility, not as the canonical contract. Consumers should use
the versioned envelope, canonical directory, metadata and copied schema.

\section{From Raw Finding to Canonical Occurrence}
\begin{figure}[H]
\centering
\begin{tikzpicture}[
 b/.style={draw=ansblue,rounded corners,fill=anslightblue,text width=28mm,
 minimum height=10mm,align=center,font=\small},
 a/.style={-{Latex[length=2mm]},thick,ansblue},node distance=7mm]
\node[b] (native) {Native provider record};
\node[b,right=of native] (finding) {Normalised finding};
\node[b,right=of finding] (identity) {Concept and identity keys};
\node[b,right=of identity] (occ) {Canonical occurrence};
\node[b,below=of occ] (evidence) {Merged provider evidence};
\node[b,left=of evidence] (status) {Semantic status and policy severity};
\node[b,left=of status] (export) {Validated projections};
\draw[a] (native)--(finding); \draw[a] (finding)--(identity);
\draw[a] (identity)--(occ); \draw[a] (occ)--(evidence);
\draw[a] (evidence)--(status); \draw[a] (status)--(export);
\end{tikzpicture}
\caption{Transformation from provider-native evidence to consumer outputs.}
\end{figure}

Raw evidence is deliberately retained. If ansible-lint and an internal rule
observe the same concept at the same sufficiently specific source location,
canonicalisation can form one occurrence with two evidence entries. The
provider raw total therefore answers ``how much did tools emit?'' whereas the
canonical total answers ``how many normalised evidence units remain?'' Neither
answer is a manual defect count.

\section{Concrete Versioned JSON Shapes}
The examples below are schematic and omit optional fields; names and nesting
match the v3 contract.

\begin{lstlisting}[language={},basicstyle=\ttfamily\footnotesize]
{
  "schema_version": "3",
  "kind": "occurrences",
  "contract": { "profile": "technical-full" },
  "analysis_health": { "status": "partial", "complete": false },
  "items": [ ... ],
  "counts": { ... },
  "artifacts": { ... },
  "policy": { ... }
}
\end{lstlisting}

\begin{lstlisting}[language={},basicstyle=\ttfamily\footnotesize]
{
  "occurrence_id": "<stable canonical id>",
  "concept_id": "unguarded_variable",
  "file": "roles/web/tasks/main.yml",
  "line": 42,
  "severity": "warning",
  "scan_mode": "snapshot",
  "semantic_subject": "service_port",
  "semantic_key": "<file-aware semantic identity>",
  "detectors": ["ansibleguard"],
  "source_population": "<manifest identity>",
  "analysis_population": "<profile/provider identity>"
}
\end{lstlisting}

The status/evidence extension records why a variable remains uncertain or
unresolved. Sensitive evidence is redacted before serialisation. Provider
evidence preserves native rule ID, finding ID, fingerprint and message so that
normalisation never destroys auditability.

\section{CSV Outputs}
\subsection{Per-occurrence CSV}
The canonical occurrence CSV is convenient for sampling and pivot tables. One
row represents one canonical occurrence. Identity/location columns identify
the unit; concept/category/severity/status columns describe it; detector and
provider fields retain origin; source/analysis population fields protect
comparisons; semantic subject and fingerprints explain matching.

\subsection{All-repositories summary CSV}
The study summary has one row per repository plus the final \code{__ALL__} row.
It contains repository/SHA/snapshot status, manifest and selected denominators,
processed and structurally analysed files, diagnostics, provider health, raw
findings, canonical totals, semantic states, symbol counts, artifacts,
categories, severities, concepts and exact rules. Wide prefixes are systematic:

\begin{longtable}{p{5.2cm}p{9cm}}
\toprule Prefix/example & Meaning \\ \midrule\endhead
\code{manifest_*} & Population before the selected production-Ansible filter:
file class, artifact/document and parse-state counts. \\
\code{selected_*} & In-scope file, LOC, artifact, document and parse counts. \\
\code{scan_*} & What was parsed or structurally analysed during execution. \\
\code{provider_*} & Raw contribution from each provider. \\
\code{finding_category_*} & Canonical counts grouped by category. \\
\code{finding_concept_*} & Canonical counts grouped by normalised concept. \\
\code{finding_rule_*} & Counts grouped by exact detector rule. \\
\code{confirmed_finding_concept_*} & Concept projection restricted to the
CONFIRMED state; it is not a human TP column. \\
\code{finding_semantic_status_*} & CONFIRMED, UNCERTAIN, EXPOSED,
\code{UNRESOLVED} and \code{ORDERED_CONFLICT} counts. \\
\code{diagnostic_stage_*}, \code{diagnostic_status_*} & Parser/provider stage
and health/failure aggregation. \\
\code{symbol_kind_*} & Definition provenance such as defaults, vars,
inventory, register, loop variable, role parameter or \code{set_fact}. \\
\code{severity_*} & Error, warning and note totals. \\
\bottomrule
\end{longtable}

\section{SARIF, Text, and Markdown}
SARIF is intended for GitHub Code Scanning and IDE-compatible consumers. It
uses repository-relative artifact URIs, exact regions, normalised rule IDs and
stable partial fingerprints. In PR mode it should communicate introduced
evidence without presenting unchanged repository debt as new work.

Text is optimized for terminal inspection and includes location, severity,
rule/concept and message. Markdown is optimized for PR descriptions, reports
and human review. JSON remains the authoritative machine interchange, while
CSV is the most practical research aggregation format. A user should not parse
text output when JSON or CSV is available.

\section{How Consumers Must Interpret Zero}
A zero has meaning only when the relevant population and provider are complete.
Examples:
\begin{itemize}
  \item zero GitLeaks findings with GitLeaks \code{completed} means no configured
  pattern matched the scanned scope;
  \item zero GitLeaks findings with GitLeaks \code{unavailable} means nothing
  about secrets;
  \item zero findings in a changed-file partial run means no selected file-local
  rule fired, not that the repository is clean;
  \item zero occurrences after canonicalisation does not erase raw diagnostics;
  \item an absent category column and a present category with value zero should
  not be conflated by downstream scripts.
\end{itemize}

This no-clean-on-failure rule is why provider health, diagnostics, profile,
scan mode, manifest identity and schema version are part of the output rather
than optional log messages.
